\documentclass[11pt]{article}
\usepackage[a4paper,margin=22mm]{geometry}
\usepackage[no-math]{fontspec}
\setmainfont{DejaVu Serif}
\usepackage{amsmath,amssymb,amsthm,mathtools,graphicx,xcolor,hyperref,xurl,xspace}
\setlength{\emergencystretch}{4em}
\SetMathAlphabet{\mathbf}{normal}{OT1}{cmr}{bx}{n}
\SetMathAlphabet{\mathbf}{bold}{OT1}{cmr}{bx}{n}
\newtheorem{theo}{Theorem}[section]
\newtheorem{lemm}[theo]{Lemma}
\newtheorem{lem}[theo]{Lemma}
\newtheorem{lemma}[theo]{Lemma}
\newtheorem{prop}[theo]{Proposition}
\newtheorem{coro}[theo]{Corollary}
\newtheorem{corol}[theo]{Corollary}
\newtheorem{conj}[theo]{Conjecture}
\newtheorem{Proposition}[theo]{Proposition}
\newtheorem{theoreme}[theo]{Theorem}
\theoremstyle{definition}
\newtheorem{defi}[theo]{Definition}
\newtheorem{exam}[theo]{Example}
\newtheorem{exem}[theo]{Example}
\newtheorem{rema}[theo]{Remark}
\newtheorem{nota}[theo]{Notation}
\newtheorem*{rema*}{Remark}
\newtheorem*{defi*}{Definition}
\newtheorem*{theo*}{Theorem}
\newtheorem*{lemm*}{Lemma}
\newtheorem*{prop*}{Proposition}
\newcommand{\enoncename}{}
\newtheorem{corpusenonce}[theo]{\enoncename}
\newenvironment{enonce}[1]{\renewcommand{\enoncename}{#1}\begin{corpusenonce}}{\end{corpusenonce}}
\newenvironment{enonce*}[1]{\par\medskip\noindent\textbf{#1.}\itshape}{\par\medskip}
\newenvironment{proofc}{\begin{proof}}{\end{proof}}
\newcommand{\Mc}{\mathcal{M}}
\newcommand{\pointir}{.\ ---\ }
\providecommand{\mathof}[1]{\mathrm{#1}}

\numberwithin{equation}{section}
\numberwithin{figure}{section}
\numberwithin{table}{section}
\def\endappendix{}
%~Mouliné par MaN_auto v.0.23.0 2020-07-08 12:32:01


\usepackage{fancybox}
\usepackage{framed}
%\usepackage{enumerate}



 	


\newcommand{\RR}{\mathbb{R}}
\newcommand{\NN}{\mathbb{N}}
\newcommand{\CC}{\mathbb{C}}
\newcommand{\TT}{\mathbb{T}}
\renewcommand{\SS}{\mathbb{S}}

\newcommand{\Cc}{\mathcal{C}}
\newcommand{\Ac}{\mathcal{A}}
\newcommand{\Bc}{\mathcal{B}}
\newcommand{\Fc}{\mathcal{F}}
\newcommand{\Lc}{\mathcal{L}}
\newcommand{\Pc}{\mathcal{P}}
\newcommand{\Qc}{\mathcal{Q}}
\newcommand{\Oc}{\mathcal{O}}
\newcommand{\drm}{\mathrm{d}}

\DeclarePairedDelimiter{\VERT}{|\mkern -1.6mu|\mkern -1.6mu|}{|\mkern -1.6mu|\mkern -1.6mu|}

%\newcommand{\Kc}{\mathcal{K}}
%\newcommand{\Uc}{\mathcal{U}}
%\newcommand{\Vc}{\mathcal{V}}
%\newcommand{\Wc}{\mathcal{W}}
%\newcommand{\Nc}{\mathcal{N}}


\newcommand{\grad}{\nabla}
\renewcommand{\div}{\operatorname{div}}
\newcommand{\cnvh}{convhull}

%{c}
\newcommand{\vect}[2]{\left(
\begin{matrix} #1 \\ #2
\end{matrix}\right)}
\newcommand{\trans}{{^{\intercal}}}

\newcounter{stepcount}
%%\renewcommand{\thestepcount}{\alph{stepcount}}
\newenvironment{step}{\refstepcounter{stepcount}\begin{proof}[Step~\thestepcount]}{\let\qed\relax\end{proof}}



%\newdimen\texpscorrection
%\texpscorrection=0truecm
%\newdimen\figcenter
%\newcommand*\figurewithtex [1] #2 #3 #4 #5\\\null \goodbreak\figcenter=\hsize\relax
%\advance\figcenter by -#4truecm
%\divide\figcenter by 2
%
%\begin{figure}[hbt]
%\vskip #3truecm\noindent\hskip\figcenter
%\special{psfile=#1}{\hskip\texpscorrection\failedinput{\baselineskip=0.8\baselineskip
%\noindent \vbox{\noindent {\footnotesize #5}}\end{figure}}}
%\newcommand*\point[1] #2 #3 {\rlap{\kern #1 truecm\raise #2 truecm \hbox{#3}}}




%%%%%%%%%%%%%%%%%%%%%%%%%%%%%%%%%%%%%%%%%%%%%%%%%


\graphicspath{{./figures/}}

\newcommand*{\mk}{\mkern -1mu}
\newcommand*{\Mk}{\mkern -2mu}
\newcommand*{\mK}{\mkern 1mu}
\newcommand*{\MK}{\mkern 2mu}

\hypersetup{urlcolor=purple, linkcolor=blue, citecolor=red}

\newcommand*{\romanenumi}{\renewcommand*{\theenumi}{\roman{enumi}}}
\newcommand*{\Romanenumi}{\renewcommand*{\theenumi}{\Roman{enumi}}}
\newcommand*{\alphenumi}{\renewcommand*{\theenumi}{\alph{enumi}}}
\newcommand*{\Alphenumi}{\renewcommand*{\theenumi}{\Alph{enumi}}}
\let\oldtilde\tilde
\renewcommand*{\tilde}[1]{\mathchoice{\widetilde{#1}}{\widetilde{#1}}{\oldtilde{#1}}{\oldtilde{#1}}}\let\oldexists\exists
\renewcommand*{\exists}{\mathrel{\oldexists}}
\let\oldforall\forall
\renewcommand*{\forall}{\mathrel{\oldforall}}

%%%%%%%%%%%%%%%%%%%%%%%%%%%%%%%%%%%%%%%%%%%%%%%%
\begin{document}
\section*{\raggedright 2638: Semilinear wave stabilization on compact connected surfaces with integrable semi-uniform linear decay and a boundary-free pseudoconvex foliation sweeping from the damping region}
\noindent\textbf{Author:} Romain Joly; Camille Laurent\par\noindent\textbf{Source:} Decay of semilinear damped wave equations: cases without geometric control condition. Annales Henri Lebesgue 3 (2020), publisher version, DOI 10.5802/ahl.60. \url{https://ahl.centre-mersenne.org/item/AHL_2020__3__1241_0/}
\par\noindent\textbf{License:} CC BY 4.0, \url{https://creativecommons.org/licenses/by/4.0/}. The original publisher PDF has the explicit license notice. See the saved source/license records.
\par\noindent\textbf{Editorial material note (not author text):} The selected problem is Theorem 1.2 in its boundary-free-foliation branch, under the exact two-dimensional compact connected manifold, nonnegative damping/nonlinearity sign, polynomial-growth and integrable semi-uniform linear-decay hypotheses. The explicit start-in-damping and full-coverage-up-to-null-set sweep conditions of Theorem 6.5 are imposed. Author introductory setup, Sections 2--3, the complete boundary-free part of Section 6 and final proof paragraph retain all necessary local asymptotic compactness, energy reduction, bowed-surface and first-contact continuation arguments. One original foliation illustration remains a figure-only asset with editable caption. Boundary-meeting, analytic-nonlinearity, quantitative-rate and example claims are unselected; references to their printed locators are preserved without importing those cores. All proof text and formulas below are editable author TeX. Mathematical wording, language and proof order are retained. Only the article class, title matter and bibliography presentation are adapted for independent compilation; no proof has been generated or repaired.
\par\noindent\textbf{Editorial difficulty:} H2: The proof derives nonlinear asymptotic compactness despite loss of uniform linear decay, identifies complete zero-dissipation limit trajectories, and eliminates them through a time-bowed pseudoconvex first-contact continuation argument. This nonlinear stabilization and geometric continuation framework exceeds ordinary qualifying exams.
\medskip\hrule\medskip
\noindent\textbf{Author statement, complete proof and local supporting text follow in source order.}
\section{Introduction}

We consider the semilinear damped wave equation
\begin{equation}\label{eq}
\begin{cases}
\partial_{tt}^2 u + \gamma(x)\partial_t u = \Delta u - \alpha u - f(x,u)\:&(x,t)\in\Omega\times (0,+\infty)\\
u_{|\partial\Omega}(x,t)=0&(x,t)\in\partial\Omega\times (0,+\infty)\\
\left(u(\cdot,t=0),\partial_t u(\cdot,t=0)\right)=U_0=(u_0,u_1)& \mkern43mu\in H^1_0(\Omega)\times L^2(\Omega)
\end{cases}
\end{equation}
in the following general framework:
\begin{enumerate}\romanenumi
\item \label{listintro1.1}the domain $\Omega$ is a two-dimensional smooth compact and connected manifold with or without smooth boundary. If $\Omega$ is not flat, $\Delta$ has to be taken as Beltrami Laplacian operator.
\item \label{listintro1.2}the constant $\alpha\geq 0$ is a non-negative constant. We require that $\alpha>0$ in the case without boundary to ensure that $\Delta-\alpha Id$ is a negative definite self-adjoint operator.
\item \label{listintro1.3}the damping $\gamma\in L^\infty(\Omega,\RR_+)$ is a bounded function with non-negative values. Since we want to consider a damped equation, we will assume that $\gamma$ does not vanish everywhere.
\item \label{listintro1.4}the non-linearity $f\in\Cc^1(\overline\Omega\times\RR,\RR)$ is of polynomial type in the sense that there exists a constant $C$ and a power $p\geq 1$ such that for all $(x,u)\in \overline\Omega\times\RR$,
\begin{equation}\label{hyp-f-estim}
|f(x,u)|+|\grad_xf(x,u)|\leq C(1+|u|)^p\:\text{ and }\:|f'_u(x,u)|\leq C(1+|u|)^{p-1}\,.
\end{equation}
Moreover, in most of this paper, we will be interested in the stabilization problem and we will also assume that
\begin{equation}\label{hyp-f-sign}
\forall (x,u)\in \overline\Omega\times\RR\,,\:f(x,u)u\geq 0\,.
\end{equation}
\end{enumerate}


We introduce the space $X=H^1_0(\Omega)\times L^2(\Omega)$ and the operator $A$ defined by
\[
D(A)=\left(H^2(\Omega)\cap H^1_0(\Omega)\right)\times H^1_0(\Omega)\quad
A =
\begin{pmatrix}
0 & Id \\
\Delta-\alpha Id & -\gamma(x)
\end{pmatrix}\,.
\]
In this paper, we are interested in the cases where the linear semigroup $e^{At}$ has no uniform decay, that is that $\|e^{At}\|_{\Lc(X)}$ does no converge to zero. We only assume a semi-uniform decay, but sufficiently fast in the following sense.

\begingroup
\begin{enumerate}\romanenumi
\setcounter{enumi}{4}
\item \label{listintro1.5}There exists a function $h(t)$ such that
\begin{equation}\label{hyp-decay-1}
\forall U_0\in D(A)\,,\,\left\|e^{At}U_0\right\|_X \, \leq \, h(t)\,\|U_0\|_{D(A)}
\end{equation}
and there is $\sigma_h\in(0,1]$ such that
\begin{equation}\label{hyp-decay-2}
\lim_{t\rightarrow \infty}h(t)=0\;\text{ and }\,
\forall \sigma\in[0,\sigma_h]\,,\:\int_0^\infty h(t)^{1-\sigma}\,\drm t\,<\,\infty\,.
\end{equation}
\end{enumerate}
\endgroup
Condition~\eqref{hyp-decay-2} requires a decay rate fast enough to be integrable. Roughly speaking, this article shows that this condition, together with a suitable unique continuation property, are sufficient to obtain a stabilization of the semilinear equation. The relevant unique continuation property is explained in Proposition~\ref{prop-UCP-implique-th} below. We present two general results where it can be obtained.

Our first result concerns analytic nonlinearities and smooth dampings.

\begin{theo}\label{th1}
Consider the damped wave equation~\eqref{eq} in the framework of Assumptions~\eqref{listintro1.1}-\eqref{listintro1.5}. Assume in addition that:
\begin{enumerate}\alphenumi
\item \label{listtheo1.1.1}the function $(x,u)\mapsto f(x,u)$ is smooth and analytic with respect to $u$.
\item \label{listtheo1.1.2}the damping $\gamma$ is of class $\Cc^1$ or at least that there exists $\tilde\gamma\in\Cc^1(\overline\Omega,\RR_+)$ such that~\eqref{listintro1.5} holds with $\gamma$ replaced by $\tilde\gamma$ and such that the support of $\tilde\gamma$ is contained in the support of $\gamma$.
\item \label{listtheo1.1.3}the power $p$ of $f$ in~\eqref{hyp-f-estim} and the decay rate $h(t)$ of the semigroup in~\eqref{hyp-decay-1} satisfy $h(t)=\Oc(t^{-\beta})$ with $\beta>2p$.
\end{enumerate}
Then, any solution $u$ of~\eqref{eq} satisfies
\[
\|(u,\partial_t u)(t)\|_{H^1_0\times L^2}\,\xrightarrow[\:t\longrightarrow
+\infty\:]{}\,0\,.
\]
Moreover, for any $R$ and $\sigma>0$, there exists $h_{R,\sigma}(t)$ which goes to zero when $t$ goes to $+\infty$ such that the following stabilization hold. For any $U_0\in H^{1+\sigma}_0(\Omega)\times H^\sigma(\Omega)$, if $u$ is the solution of~\eqref{eq}, then
\[
\|(u_0,u_1)\|_{H^{1+\sigma}\times H^\sigma}\leq
R\:\Longrightarrow\:\|(u,\partial_t u)(t)\|_{H^1_0\times L^2} \leq
h_{R,\sigma}(t)\,\xrightarrow[\:t\longrightarrow +\infty\:]{}\,0\,.
\]
\end{theo}

Our assumptions~\eqref{listintro1.5} and~\eqref{listtheo1.1.3} on the decay of the linear semigroup may seem strong. They are satisfied in the cases where the set of trapped geodesics, the ones which do not meet the support of the damping, is small and hyperbolic in some sense. Several geometries satisfying~\eqref{listintro1.5} and~\eqref{listtheo1.1.3} have been studied in the literature, see the concrete examples of Figure~1.1 and the references therein. Notice in particular that the example of domain with holes is particularly relevant for applications where we want to stabilize a nonlinear material with holes by adding a damping or a control in the external part. There is a huge literature about the damped wave equation and the purpose of the examples presented here is mainly to illustrate our theorem to non specialists. Moreover, the subject is growing quickly, giving more and more examples of geometries where we understand the effect of the damping and where we may be able to apply our results. We do not pretend to exhaustivity and refer to the bibliography of the more recent~\cite{Lin} for instance.

In some cases, the unique continuation property required in Proposition~\ref{prop-UCP-implique-th} can be obtained without considering analytic nonlinearities or conditions on the growth of $f$ as Hypothesis c) of Theorem~\ref{th1}. Instead, we require a particular geometry, which will be introduced more precisely in Section~\ref{section-UCP}.
\goodbreak

\begin{theo}\label{th2}
Consider the damped wave equation~\eqref{eq} in the framework of Assumptions~\eqref{listintro1.1}-\eqref{listintro1.5}. Assume in addition that:
\begin{enumerate}\alphenumi
\item \label{listtheo1.2.1}the function $(x,u)\mapsto f(x,u)$ is of class $\Cc^1\left(\overline{\Omega}\times \RR,\RR\right)$.
\item \label{listtheo1.2.2}there exists a pseudo-convex foliation of $\Omega$ in the sense of Definitions~\ref{defi-foliation1} or~6.6.
\end{enumerate}
Then, the conclusions of Theorem~\ref{th1} hold.
\end{theo}

This result can be applied in several situations of Figure~1.1: the ``disk with two holes'', the ``peanut of rotation'' and the ``open book''. In these cases, the stabilization holds for any natural nonlinearity.

We expect that the decay rate $h_{R,\sigma}(t)$ is related to the linear decay rate $h(t)$ of Assumption~\eqref{listintro1.5}. We are able to obtain this link for the typical decays of the examples of Figure~1.1.

\begin{prop}\label{prop1}
Consider a situation where the stabilization stated in Theorems~\ref{th1} or~\ref{th2} holds. Then,
\begin{itemize}
\item if the decay rate of Assumption~\eqref{listintro1.5} satisfies $h(t)=\Oc({t^{-\alpha}})$ with $\alpha>1$, then the nonlinear equation admits a decay of the type $h_{R,\sigma}(t)= \Oc (t^{-\sigma\alpha})$.
\item if the decay rate of Assumption~\eqref{listintro1.5} satisfies $h(t)=\Oc \left(e^{-a {t^{1/\beta}}}\right)$ with $a>0$ and $\beta>0$, then the nonlinear equation admits a decay of the type $h_{R,\sigma}(t)\linebreak= \Oc \left(e^{-b\sigma t^{1/(\beta+1)}}\right)$ for some $b>0$.
\end{itemize}
\end{prop}

Notice that this result is purely local in the sense that the decay rate is obtained when the solution is close enough to $0$. Our proofs do not provide an explicit estimate of the time needed to enter this small neighborhood of $0$. Also notice that the loss in the power of the second case of Proposition~\ref{prop1} is due to an abstract setting: in the concrete examples, we may avoid this loss, see the remark below Lemma~4.2 and the concrete applications to the examples of Figure~1.1.

\section{Notations and basic facts}\label{section-basic}
We use the notations of Equation~\eqref{eq}, of Assumptions~\eqref{listintro1.1}-\eqref{listintro1.5} and of the introduction. In particular, we recall that $X=H^1_0(\Omega)\times L^2(\Omega)$ and
\[
D(A)=\left(H^2(\Omega)\cap H^1_0(\Omega)\right)\times H^1_0(\Omega)\,
A =
\begin{pmatrix}
0 & Id \\ \Delta-\alpha Id & -\gamma(x)
\end{pmatrix}\,.
\]
The operator $A$ is the classical linear damped wave operator corresponding to the linear part of~\eqref{eq}. Due to Lumer--Phillips theorem, we know that this operator generates a linear semigroup of contractions $e^{At}$ on $X$ and on $D(A)$ and that
\[
\forall t\geq 0\,,\,\left\|e^{At}\right\|_{\Lc(X)}\leq 1\quad\text{and}\quad\left\|e^{At}\right\|_{\Lc(D(A))}\leq 1\,.
\]
Notice that the second estimate is a direct consequence of the commutation of $A$ and $e^{At}$ and does not require any regularity on $\gamma$.

For any $\sigma\in [0,1]$, we set
\[
X^\sigma=\left(H^{1+\sigma}(\Omega)\cap H^1_0(\Omega)\right)\times
H^\sigma_0(\Omega)\,.
\]
Thus $X^0=X$ and $X^1=D(A)$ and $X^\sigma$ is an interpolation space between $X^0$ and $X^1$. In particular, by interpolation, $e^{At}$ is defined in $X^\sigma$ and we have
\begin{equation}\label{eq-bound-semigroup}
\forall \sigma\in [0,1]\,,\:\forall t\geq 0\,,\:
\left\|e^{At}\right\|_{\Lc(X^\sigma)}\leq 1\,.
\end{equation}

We set $F\in\Cc^0(X)$ to be the function
\begin{equation}\label{def-F}
F:U=\vect uv \in X\,\longmapsto\,\vect 0 {-f(\cdot,u)} \in X\,.
\end{equation}
Notice that, if $\Omega$ is two-dimensional, $H^1_0(\Omega) \hookrightarrow L^{p}(\Omega)$ for any $p\in[1,+\infty)$ and if $\Omega$ is three-dimensional $H^1_0(\Omega) \hookrightarrow L^6(\Omega)$. Thus, $f(u)$ is well defined in $L^2(\Omega)$ due to Assumption~\eqref{hyp-f-estim} if dim$(\Omega)=2$ or if dim$(\Omega)=3$ and $p\leq 3$. Moreover, for any $R$, $u$ and $v$ with $\|u\|_{H^1}\leq R$ and $\|v\|_{H^1}\leq R$, we have
\[
\|f(\cdot,u)-f(\cdot,v)\|_{L^2}=\left\|(u-v)\int_0^1
f'_u(\cdot,u+s(u-v))ds\right\|_{L^2}
\leq C(R)\|u-v\|_{H^1}
\]
and so $F$ is lipschitzian on the bounded sets of $X$. As a consequence, the damped wave equation~\eqref{eq} is well posed in $X$ and admits local solutions if dim$(\Omega)=2$ or if dim$(\Omega)=3$ and $p\leq 3$.

With the above notation, our main equation writes
\begin{equation}\label{eq2}
\partial_t U=AU + F(U)\quad U(t=0)=U_0\in X\,.
\end{equation}
In particular, Duhamel's formula yields
\[
U(t)=e^{At}U_0 + \int_0^t e^{A(t-s)}F(U(s))\,\drm s\,.
\]

We introduce the potential
\[
V(x,u)=\int_0^u f(x,s)ds\,.
\]
Due to~\eqref{hyp-f-estim} and the above arguments, $V(\cdot,u)$ defines a Lipschitz function from the bounded sets of $H^1_0(\Omega)$ into $L^1(\Omega)$. The classical energy associated to~\eqref{eq} is defined along a trajectory $U=(u,\partial_t u)$ as
\[
E(U)= \int_\Omega \frac12\left(|\grad u|^2+\alpha |u|^2+|\partial_t u|^2\right) +
V(x,u)\,.
\]
The damping effect appears by the computation
\begin{equation}\label{eq-deriv-E}
\partial_t E(U(t))= - \int_\Omega \gamma(x)|\partial_t u |^2\,.
\end{equation}
In particular, the energy $E$ is non-increasing along the trajectories. Moreover, the sign assumption~\eqref{hyp-f-sign} yields that $V(x,u)\geq 0$. Thus, we have that $E(U)\geq C \|U\|_X^2$ and that $E(t)$, $t\geq 0$, is bounded on the bounded sets of $X$. All together, the above properties show that for any $U_0\in X$, the solution $U=(u,\partial_t u)$ of~\eqref{eq} is defined for all non-negative times and remains in a bounded set of $X$, which only depends on $\|U_0\|_X$.

A fundamental question of this paper concerns the solution for which the energy is constant: are they equilibrium points or may they be moving trajectories? At least, the answer is known for the linear equation, see~\cite{Dafermos} and also~\cite{Haraux}.

\begin{theo}\label{th-stab-linear}
\emph{Dafermos (1978).}
Assume that the damping $\gamma\geq 0$ does not vanish everywhere. Then, for any $U_0\in X$, we have
\[
e^{At}U_0\xrightarrow[t\longrightarrow +\infty]{} 0\text{ in }X\,.
\]
\end{theo}



\section{Asymptotic compactness and reduction to a unique continuation problem}\label{section-asymp-compact}

In this section, we assume a fast enough semi-uniform linear decay as described by~\eqref{hyp-decay-1} and~\eqref{hyp-decay-2}. We first notice that, by linear interpolation, we have the following result.

\begin{prop}\label{prop-decay-interpole}
For any $\sigma_1,\sigma_2$ such that $0\leq \sigma_2<\sigma_1 \leq 1$, the linear semigroup is well defined from $X^{\sigma_1}$ in $X^{\sigma_2}$ and we have
\[
\forall U_0\in X^{\sigma_1}\,,\,\left\|e^{At}U_0\right\|_{X^{\sigma_2}} \leq
h(t)^{\sigma_1-\sigma_2} \|U_0\|_{X^{\sigma_1}}\,.
\]
\end{prop}
\begin{proof}
We interpolate the estimates~\eqref{eq-bound-semigroup} for $\sigma=1$ and~\eqref{hyp-decay-1} with respective weights $(\sigma_2/\sigma_1,1-\sigma_2/\sigma_1)$ and obtain
\[
\forall U_0\in D(A)\,,\,\left\|e^{At}U_0\right\|_{X^{\sigma_2/\sigma_1}} \leq
h(t)^{1-\sigma_2/\sigma_1}\|U_0\|_{D(A)}\,.
\]
It remains to interpolated the above estimate and~\eqref{eq-bound-semigroup} for $\sigma=0$ with respective weights $(\sigma_1,1-\sigma_1)$.
\end{proof}



We also need some regularity properties for $F$. The following properties depend on Sobolev embeddings and so of the dimension $d$ of $\Omega$. For $d=2$, which is the case in our examples, the properties are general. For $d=3$, they are more restrictive but they are shown in the same way. We choose to also consider this case in our paper for possible later uses.

\begin{prop}\label{prop-compact-F}
Assume that dim$(\Omega)=2$. Then for any $\sigma\in [0,1)$, the function $F$ maps any bounded set $\Bc$ of $X=H^1_0(\Omega)\times L^2(\Omega)$ in a bounded set $F(\Bc)$ contained in $X^\sigma= (H^{1+\sigma}(\Omega)\cap H^1_0(\Omega))\times H^\sigma_0(\Omega)$. Moreover, $F(\Bc)$ has compact enclosure in $X^\sigma$.

If dim$(\Omega)=3$ and if~\eqref{hyp-f-estim} holds for some $p\in [0,3)$, then the same properties hold for $\sigma\in [0,(3-p)/2)$.
\end{prop}
\begin{proof}
Assume that dim $(\Omega)=2$. First notice that we only have to show that $f(\cdot,u)$ is compactly bounded in $H^\sigma_0(\Omega)$ since the first component of $F$ is zero. Also notice that $f(x,0)=0$ due to the sign Assumption~\eqref{hyp-f-sign}, thus the Dirichlet boundary condition possibly contained in $H^\sigma_0(\Omega)$ will be fulfilled by $f(x,u)$ if $u\in H^1_0(\Omega)$. Due to the Sobolev embeddings, and since $\Omega$ is compact, it is sufficient to show that $F(\Bc)$ is bounded in $W^{1,q}(\Omega)$ for all $q\in [1,2)$ to obtain compactness in $H^\sigma(\Omega)$ for any $\sigma\in [0,1)$. Since $f$ is of polynomial type due to~\eqref{hyp-f-estim}, we know that $f(x,u)$ is bounded in $L^q(\Omega)$ for any $q\in [1,2)$. On the other hand, using~\eqref{hyp-f-estim}, we have
\begin{align*}
\|\grad(f(x,u))\|_{L^q} & \leq \|\grad_xf(x,u)\|_{L^q} + \left\|f'_u(x,u)\grad u\right\|_{L^q}\\
& \leq C \|(1+|u|)^p \|_{L^q} + C \left\|(1+|u|)^{p-1} \grad u\right\|_{L^q}\\
& \leq C\left(1+\|u\|_{L^{pq}}^p + \|\grad u\|_{L^2} \|u\|^{p-1}_{L^r}\right)
\end{align*}
with $r=(p-1)\frac {2q}{2-q}$ defined as soon as $q<2$. This shows that $f(\cdot,u)$ belongs to $W^{1,q}(\Omega)$ for any $q\in [1,2)$ and concludes the proof for dim$(\Omega)=2$.

The case dim$(\Omega)=3$ is similar once we use the suitable Sobolev embeddings.
\end{proof}


The main results of this section are the following asymptotic compactness properties.

\begin{prop}\label{prop-asymp-compact-1}
If dim$(\Omega)=2$, set $\sigma_*=\sigma_h$. If dim$(\Omega)=3$, assume that $p\in(1,3)$ in~\eqref{hyp-f-estim} and that $\sigma_h\in ((p-1)/2,1)$ in~\eqref{hyp-decay-2} and set $\sigma_*=\sigma_h - (p-1)/2$.

Let $U(t)=(u,\partial_t u)$ where $u$ solves~\eqref{eq} and let $(t_n)$ be a sequence of times such that $t_n\rightarrow +\infty$. Then, there exist a subsequence $(t_{\varphi(n)})$ and a solution $W(t)=(w,\partial_t w)(t)$ of~\eqref{eq} defined for all $t\in\RR$, such that
\[
\forall
t\in\RR\,,\,U(t_{\varphi(n)}+t)\,\xrightarrow[\:n\longrightarrow
+\infty\;]{}\,W(t)\;\text{ in }X=H^1_0(\Omega)\times L^2(\Omega)\,.
\]
Moreover, the solution $W$ is globally bounded in $X^\sigma$ for all $\sigma\in [0,\sigma_*)$ and the energy $E(W(t))$ is constant.
\end{prop}
\begin{proof}
Assume first that dim$(\Omega)=2$. We have
\begin{equation}\label{eq-prop-asymp-compact-1}
U(t_n)=e^{At_n}U_0 + \int_0^{t_n} e^{A(t_n-s)}F(U(s))\,\drm s\,.
\end{equation}
Due to Theorem~\ref{th-stab-linear}, the term $e^{At_n}U_0$ goes to zero in $X$. Thus, it remains to show that $\int_0^{t_n} e^{A(t_n-s)}F(U(s))\,\drm s$ is a compact term in $X$. First notice that $U(s)$ is uniformly bounded for $s\geq 0$ due to the non-increasing energy $E$ (see Section~\ref{section-basic}). Due to Proposition~\ref{prop-compact-F}, $F(U(s))$ thus belongs to a bounded set of $X^{\sigma_1}$ for all $\sigma_1\in [0,1)$. By Assumption~\eqref{hyp-decay-2} and Proposition~\ref{prop-decay-interpole}, $e^{A(t_n-s)}F(U(s))$ has an integral in $[0,t_n]$ bounded in $X^{\sigma_2}$ uniformly with respect to $n$, for any $\sigma_2 \in (0,\sigma_h)$. Thus $\int_0^{t_n} e^{A(t_n-s)}F(U(s))\,\drm s$ is a compact sequence in $X^\sigma$ for any $\sigma\in [0,\sigma_2)$. As a consequence, for any $\sigma$ as close as wanted to $\sigma_h$, we may extract a subsequence $(t_{\varphi(n)})$ such that $\int_0^{t_{\varphi(n)}} e^{A(t_{\varphi(n)}-s)}F(U(s))\,\drm s$ converges to some limit $W(0)$ in $X^\sigma$. Since the linear term of~\eqref{eq-prop-asymp-compact-1} goes to zero in $X$ for $t_{\varphi(n)}\rightarrow +\infty$, $U(t_{\varphi(n)})$ converges to $W(0)\in X^\sigma$ for the norm of $X$.

Let $W(t)=(w,\partial_t w)(t)$ be the maximal solution of the damped wave equation~\eqref{eq} corresponding to the initial data $W(0)$. Let $t\in\RR$, for $n$ large enough $t_{\varphi(n)}+t\geq 0$ and thus $U(t_{\varphi(n)}+t)$ is well defined and uniformly bounded in $X$. Since our equation is well posed, the solution is continuous with respect to the initial data. Thus, since $U(t_{\varphi(n)})$ converges to $W(0)$ in $X$, we have that $U(t_{\varphi(n)}+t)$ converges to $W(t)$ for all $t$ such that $W(t)$ is well defined. But due to the uniform bound on $U(t_{\varphi(n)}+t)$, $W(t)$ is uniformly bounded and thus the solution may be extended to a global solution $W(t)$, $t\in\RR$. In addition, the $X^\sigma$-bound obtained above for $W(0)$ only depends on the $X-$bound on $U(s)$ which is uniform due to non-increase of the energy of $U(t)$. Thus, the same arguments applied to the convergence $U(t_{\varphi(n)}+t)\rightarrow W(t)$ give the same $X^\sigma$-bound for $W(t)$ for all $t\in\RR$. Finally, since the energy of $U(t)$ is non-increasing and non-negative, for any $t\in\RR$, we must have $E(U(t_n+t))-E(U(t_n))\rightarrow 0$ when $n\rightarrow 0$ (since $t_n$ goes to $+\infty$). This shows that $E(W(t))$ is constant and finishes the proof.

The case dim$(\Omega)=3$ is similar once we take into account the constraints given by Proposition~\ref{prop-compact-F}.
\end{proof}


\goodbreak

\begin{prop}\label{prop-asymp-compact-2}
If dim$(\Omega)=2$, set $\sigma_*=\sigma_h$. If dim$(\Omega)=3$, assume that $p\in(1,3)$ in~\eqref{hyp-f-estim} and that $\sigma_h\in ((p-1)/2,1)$ in~\eqref{hyp-decay-2} and set $\sigma_*=\sigma_h - (p-1)/2$.

Let $\sigma>0$ and $R>0$. Let $U_n(t)=(u_n,\partial_t u_n)$ a sequence of solutions $u_n$ of~\eqref{eq} such that $(U_n(0))\subset X^\sigma$ and $\|U_n(0)\|_{X^\sigma}\leq R$. Let $(t_n)$ be a sequence of times such that $t_n\rightarrow +\infty$ and let $\sigma'\in [0,\min(\sigma,\sigma_*))$. Then, there exist subsequences $(t_{\varphi(n)})$ and $(U_{\varphi(n)})$ and a solution $W(t)=(w,\partial_t w)(t)$ of~\eqref{eq} defined for all $t\in\RR$, such that
\[
\forall
t\in\RR\,,\,U_{\varphi(n)}(t_{\varphi(n)}+t)\,\xrightarrow[\;n\longrightarrow
+\infty\:]{}_,W(t)\;\text{ in } X^{\sigma'}\,.
\]
Moreover, the solution $W$ is globally bounded in $X^{\sigma'}$ and the energy $E(W(t))$ is constant.
\end{prop}
\begin{proof}
The arguments are similar as the ones of the above proof of Proposition~\ref{prop-asymp-compact-1}. The term $e^{At_n}U_n(0)$ goes to zero in $X^{\sigma'}$ due to Proposition~\ref{prop-decay-interpole} because $\sigma'<\sigma$. We bound the integral $\int_0^{t_n} e^{A(t_n-s)}F(U_n(s))\,\drm s$ as in the proof of Proposition~\ref{prop-asymp-compact-1}: $U_n(s)$ is uniformly bounded in $X$, so $F(U_n(s))$ is uniformly bounded in $X^\eta$ with $\eta<1$ in dimension~$2$ or $\eta<(3-p)/2$ in dimension\,$3$. Proposition~\ref{prop-decay-interpole} together with~\eqref{hyp-decay-2} implies that the integral is uniformly bounded in $X^{\sigma'}$ with $\sigma'<\sigma_*$. The compactness follows by leaving any small amount of regularity in the process. To obtain the convergence to $W(t)$ for all $t$, we use the same argument as the one of the proof of Proposition~\ref{prop-asymp-compact-1} to first show the convergence in $X$. Then, the above arguments also show the compactness of $U(t_{\varphi(n)}+t)$ in $X^{\sigma'}$ and thus the convergence to $W(t)$ also holds in $X^{\sigma'}$. The last property is the same as the ones of Proposition~\ref{prop-asymp-compact-1}.
\end{proof}

The conclusions of Theorem~\ref{th1} then follow from Propositions~\ref{prop-asymp-compact-1} and~\ref{prop-asymp-compact-2} as soon as we can prove that $W\equiv 0$ for any subsequences of any sequences $(t_n)$ and $U_n$. To this end, notice that $E(W(t))$ is constant and its derivative~\eqref{eq-deriv-E} implies that $\int_\Omega \gamma(x) |\partial_t w|^2=0$ for all time. To formulate this property as a unique continuation property, we set as usual $z=\partial_t w$ and notice that $z$ solves
\begin{equation}\label{eq-UCP}
\begin{cases}
\partial_{tt}^2 z(x,t) = \Delta z(x,t) - \alpha z(x,t) - f'_u(x,w(x,t)) z\;&(x,t)\in\Omega\times\RR\\
z_{|\partial\Omega}(x,t)=0&(x,t)\in\partial\Omega\times \RR\\
\quad z(x,t)\equiv 0 & (x,t)\in \text{support}(\gamma)\times \RR
\end{cases}
\end{equation}
If this implies $z\equiv 0$ everywhere, this means that $w(x,t)=w(x)$ is constant in time and solves
\[
\Delta w(x)-\alpha w(x)=f(x,w(x))\,.
\]
Multiplying by $w$ and integrating, we obtain
\[
\|\grad w\|^2 +\alpha \|w\|^2 = -\int_\Omega f(x,w(x))w(x)\,\drm x\,.
\]
By the sign Assumption~\eqref{hyp-f-sign}, this yields $w\equiv 0$. Thus, it only remains to study this unique continuation property.

\begin{prop}\label{prop-UCP-implique-th}
Assume that $z\equiv 0$ is the only global solution of~\eqref{eq-UCP}. Then the decay assumptions~\eqref{hyp-decay-1} and~\eqref{hyp-decay-2} imply the conclusions of Theorem~\ref{th1}.
\end{prop}


\setcounter{section}{5}
\section{Unique continuation theorems}\label{section-UCP}
As proved in Section~\ref{section-asymp-compact}, the last step to prove stabilization is the unique continuation property: if $z$ is a global solution of
\begin{equation}\label{eq-UCP2}
\begin{cases}
\partial_{tt}^2 z(x,t) = \Delta z(x,t) - \alpha z(x,t) - f'_u(x,w(x,t)) z\;&(x,t)\in\Omega\times\RR\\
z_{|\partial\Omega}(x,t)=0&(x,t)\in\partial\Omega\times \RR\\
\mkern20mu z(x,t)\equiv 0 & (x,t)\in \text{support}(\gamma)\times \RR
\end{cases}
\end{equation}
then $z\equiv 0$. Except for the example of Section~5, the property is often difficult to obtain. The purpose of this section is to gather several results yielding this property.

The first known result has been proved by Ruiz in~\cite{Ruiz}. It stated the unique continuation property in a bounded domain $\Omega \subset\RR^d$ as soon as the support of $\gamma$ contains a neighborhood of the boundary $\partial\Omega$. This result has been generalized in~\cite{K-K} (see also~\cite{Las-Tri-Zha} for Neumann boundary conditions). However, this kind of results is not relevant in this paper. Indeed, their geometric settings implies the uniform decay of the semigroup $e^{At}$ and we are interested here in cases where it is not satisfied. We need sharper results.


\subsection{Unique continuation with coefficients analytic in time}

A very general unique continuation property holds if the coefficients of a linear wave equation as~\eqref{eq-UCP2} are analytic in time. This is a consequence of local continuation results proved by H\"ormander in~\cite{Hormander_uniq_analytic} and generalized by Tataru in~\cite{Tataru_uniq_analyticJMPA} and also independently proved by Robbiano and Zuily in~\cite{Rob-Zui}. These results concern in fact a very general setting but we restrict here the statements at the case of the wave equation. The application to the wave equation and the proof that the local results yield a global one are classical and straightforward, see for example~\cite[Corollary~3.2]{RC} for the details.

\begin{theo}\label{th-RZ}
{\emph{Robbiano--Zuily, H\"ormander (1998).}}
Let $T>0$ (or $T=+\infty$) and let $b$, $c$ and $d$ be smooth coefficients. Assume moreover that $b$, $c$ and $d$ are analytic in time and that $z$ is a strong solution of
\begin{equation}\label{eq-RZ}
\begin{cases}
\partial^2_{tt} z &=\Delta z + b(x,t)\partial_t z+c(x,t).\grad z+d(x,t)z\quad(x,t)\in\Omega\times (-T,T)\\
z_{|\partial\Omega}(x,t)&= 0\mkern298mu(x,t)\in\partial\Omega\times \RR\,.
\end{cases}
\end{equation}
Let $\Oc$ be a non-empty open subset of $\Omega$ and assume that $z(x,t)=0$ in $\Oc\times (-T,T)$. Then $z(x,0)\equiv 0$ in $\Oc_T=\{x_0\in\Omega\,,\,d(x_0,\Oc)<T\}$.

As consequences if $z\equiv 0$ in $\Oc\times (-T,T)$ and $\overline
\Oc_T=\Omega$, then $z\equiv 0$ everywhere.
\end{theo}



\subsection{Unique continuation through pseudo-convex surfaces without boundary}\label{subsectUCPwithout}
If the coefficients of~\eqref{eq-RZ} are not analytic in time, the geometry of the problem is more constrained. However, it could still include cases where the geometric control condition of~\cite{BLR} does not hold and thus where the semigroup $e^{At}$ is not uniformly stable, see the examples below.

We consider here H\"ormander framework (see~\cite{Hormander} for example). The principal symbol of the differential operator of~\eqref{eq-UCP2} is of order two and writes locally
\[
p(x,t,\xi,\tau)=\xi\trans.A(x).\xi-|\tau|^2
\]
where $A(x)$ is a smooth family of positive definite symmetric matrices coding the Beltrami Laplacian operator in a local chart. Let $\psi(x,t,\xi,\tau)$ be a locally $\Cc^1-$function, we introduce the Poisson bracket
\[
H_p(\psi)=\{p,\psi\}= \grad_\xi p \grad_x \psi + \partial_\tau p \partial_t
\psi - \grad_x p \grad_\xi \psi - \partial_t p \partial_\tau \psi\,.
\]
Let $\psi$ be a smooth function defined in a neighborhood $\Oc\subset
\RR^{d+1}$ of $(x_0,t_0)$. Assume that $(\grad_x \psi,\partial_t
\psi)(x_0,t_0)\neq 0$ so that $\Sigma=\{(x,t),\psi(x,t)=0\}$ defines a smooth hypersurface near $(x_0,t_0)$.

\begin{defi}\label{defi-pseudoconvex}
The local hypersurface $\Sigma$ is said to be \emph{non-characteristic} at $(x_0,t_0)$ if
\[
p(x_0,t_0,\grad \psi(x_0,t_0),\partial_t \psi(x_0,t_0))\neq 0\,.
\]
Moreover, $\Sigma$ is said to be \emph{strongly pseudo-convex} at $(x_0,t_0)$ if for any $(\xi,\tau)\neq 0$ such that $p(x_0,t_0,\xi,\tau)=0$ and $H_p(\psi)(x_0,t_0)=0$, we have
\[
H_p^2(\psi)(x_0,t_0)>0\,.
\]
\end{defi}

Notice that the above Definition~\ref{defi-pseudoconvex} of strongly pseudo-convexity is adapted to the case of a real differential operator of order two. Thus it is perfectly adapted to the situation of this paper where the wave operator is $p=\xi\trans.A(x).\xi-|\tau|^2$. However, we emphasize that, in the general case, the assumption of pseudo-convexity is more complex, see~\cite{Hormander}.

The geometrical interpretation of Definition~\ref{defi-pseudoconvex} is as follows. First, the fact that the surface is non-characteristic says that $|\partial_t \psi|^2 \neq \grad\psi\trans.A(x).\grad \psi$. This means that the surface is not moving at the exact same speed as the sound waves.

The pseudo-convexity is slightly more involved. Consider the total Hamiltonian flow $\sigma\mapsto\varphi_\sigma$ defined by
\[
\varphi_0(x,t,\xi,\tau)=(x,t,\xi,\tau)\;\partial_\sigma
\varphi_\sigma(x,t,\xi,\tau)=(\grad_\xi p,\partial_\tau
p,-\grad_x p,-\partial_t p)(\varphi_\sigma)\,.
\]
Since $p$ is independent of $t$, $\tau$ is constant and thus $t(\sigma)=t-2\tau\sigma$, meaning that $\sigma$ is a simple new parametrization of time. Moreover, $\grad_\xi p$ and $\grad_x p$ are independent of $t$ and $\tau$. Thus, $(x,\xi)(\sigma)$ follows the geodesic flow
\[
\partial_\sigma (x,\xi)=(\grad_\xi g^*,-\grad_x g^*)(x,\xi)
\]
where $g^{*}(x,\xi)=\xi\trans.A(x).\xi$ is the symbol of the local cometric. Assume $p(x,t,\xi,\tau)=0$ at $\sigma=0$. The Hamiltonian being conserved, we always have $p(x,t,\xi,\tau)=0$ and $|\tau|^2=\xi\trans A(x)\xi$ is constant in $\sigma$: the point $x(\sigma)$ is moving along a geodesic of the metric at a speed which is of constant norm $|\tau|$ with respect to the metric. Let $h$ be a function of $(x,t,\xi,\tau)$, then the Poisson bracket $\{p,h\}$ is the derivative at $\sigma=0$ of $h(\varphi_\sigma(x,t,\xi,\tau))$. Thus
\[
\{p,h\}=\{g^*,h\} -2\tau \partial_t h=\{g^*,h\} + \partial_\sigma h
\]
where $\{g^*,h\}$ is the derivative along the geodesic $(x,\xi)(\sigma)$ of the metric starting at $x$ with speed $\xi$. Thus, the strongly pseudo-convexity condition $H_p\psi=0\Rightarrow H_p^2 \psi>0$ means that if a geodesic of the surface is tangent to $\Sigma$ in the space-time sense, then it must be contained in a non-degenerated sense in the half-space $\psi(x,t)>0$ for $t\neq0$. Finally notice that if $\psi$ does not depend on time $t$ (as in Definitions~\ref{defi-foliation1} and~6.6), then the strongly pseudo-convexity is a classical strong convexity: if a classical geodesic of the metric $g$ is tangent to the surface $\psi(x)=0$, it must be contained in the half-space $\psi(x)>0$.

\cite[Theorem~28.4.3]{Hormander} mainly comes from~\cite{LR} and is stated as follows.

\begin{theo}\label{th-hormander}
\emph{Lerner and Robbiano (1985), H\"ormander.}
Let $\Oc$ be a small open neighborhood of a point $(x_0,t_0)$ in $\RR^d\times
\RR$ and let $A(x)$ be a smooth family of positive definite symmetric matrices. Let $b$, $c$ and $d$ be bounded coefficients. Assume that $z$ is a mild solution of
\begin{equation}\label{eq-LR}
\partial^2_{tt} z=\div A(x)\grad z + b(x,t)\partial_t z+c(x,t).\grad z+d(x,t)z\;(x,t)\in\Oc\,.
\end{equation}
Let $\Sigma=\{(x,t),\psi(x,t)=0\}$ be a smooth surface containing $(x_0,t_0)$ which is non-characteristic and strongly pseudo-convex in the sense of Definition~\ref{defi-pseudoconvex}.

Then, if $u(x,t)=0$ for all $(x,t)\in\Oc$ such that $\psi(x,t)\geq 0$, we have $u(x,t)\equiv 0$ in a neighborhood of $(x_0,t_0)$.
\end{theo}

Theorem~\ref{th-hormander} states a local unique continuation property through pseudo-convex surfaces. To use it, it is more convenient to have a global version. This kind of global foliation has already been introduced in~\cite{Stefanov-Uhlmann} by Stefanov and Uhlmann.

\begin{defi}\label{defi-foliation1}
A family of surfaces $(\Sigma_\lambda)_{\lambda\in[0,1)}$ is an \emph{oriented pseudo-convex foliation without boundary} in a compact manifold $\Omega$ if:
\begin{enumerate}\romanenumi
\item \label{listdef6.4.1}the family of surfaces is smooth in the sense that it is locally described as level sets $\{x,\psi_\lambda(x)=0\}$ where $(x,\lambda)\mapsto\psi_\lambda(x)$ is a local smooth function with $\grad_x\psi_\lambda\neq 0$.
\item \label{listdef6.4.2}each surface is globally oriented in the sense that there exist disjoint sets $\Sigma_\lambda^\pm$ such that locally $\{x\in\Omega,\pm\psi_\lambda(x)> 0\}\subset \Sigma_\lambda^\pm$ and such that $\Omega=\Sigma_\lambda^- \cup
\Sigma_\lambda \cup \Sigma_\lambda^+$.
\item \label{listdef6.4.3}for each $\lambda$, $(x,t)\mapsto \psi_\lambda(x)$ is pseudo-convex in the sense of Definition~\ref{defi-pseudoconvex} as a function independent of $t$. Equivalently, $\Sigma_\lambda^-$ is locally strictly convex in a neighborhood of its boundary $\Sigma_\lambda$ for the metric $g$: for each $x\in\Sigma_\lambda$, a geodesic through $x$ which is tangent at $\Sigma_\lambda$ is locally included in $\Sigma_\lambda^+$, $x$ excepted.
\item \label{listdef6.4.4}the surfaces $\Sigma_\lambda$ are compact and have no boundary or equivalently do not meet $\partial\Omega$.
\end{enumerate}
\end{defi}

A typical example of such oriented pseudo-convex foliation without boundary is given in Figure~\ref{fig-foliation-1}.

\begin{figure}[ht]
\centering\includegraphics[width=171pt]{assets/boundary-free-foliation.png} 
\caption{A oriented pseudo-convex foliation without boundary in the sphere $\SS^2$. The surfaces $\Sigma_\lambda$ forms a smooth family of vertical circles inside a hemisphere and $\Sigma_\lambda^+$ and $\Sigma_\lambda^-$ are respectively the large and the small spherical caps. The geodesics beings the great circles, the ones which are tangent to a surface $\Sigma_\lambda$ stay inside~$\Sigma_\lambda^+$.\label{fig-foliation-1}
}
\end{figure}

By a classical argument, we may state a global version of Theorem~\ref{th-hormander} as follows.

\begin{theo}\label{th-hormander-global}
Let $\Omega$ be a smooth compact manifold (with or without boundaries) and let $\omega\subset\Omega$ be an open set. Assume that there exists an oriented pseudo-convex foliation without boundary $(\Sigma_\lambda)_{\lambda\in[0,1)}$ of $\Omega$ in the sense of Definition~\ref{defi-foliation1}. Also assume that $\Sigma_0^+\subset \omega$ and $\omega\cup \left(\bigcup_{\lambda\in[0,1)}
\Sigma_\lambda^+\right)$ covers $\Omega$ up to a set of zero measure.

Let $b$, $c$ and $d$ be bounded coefficients. Assume that $z$ is a global mild solution of
\[
\begin{cases}
\partial^2_{tt} z &=\Delta z + b(x,t)\partial_t z+c(x,t).\grad z+d(x,t)z\quad(x,t)\in\Omega\times \RR\\
z_{|\omega}(x,t)&= 0\mkern297mu(x,t)\in\omega\times\RR
\end{cases}
\]
with any suitable boundary conditions on $\partial\Omega$ such that the wave equation is well-posed and where $\Delta$ is the Laplace--Beltrami operator related to $\Omega$.

Then $z\equiv 0$ everywhere.
\end{theo}


\begin{proof}
We will show that $z(\cdot,t=0)$ vanishes in $\omega\cup\Sigma^+_\lambda$, for all $\lambda\in [0,1)$, which shows that $z(t=0)\equiv 0$ and thus that $z\equiv 0$ due to the uniqueness properties of the wave equation. By assumption, $z\equiv 0$ in $\Sigma_0^+\subset \omega$. Let $\lambda_0\in (0,1)$ and let $h_{\alpha,T}(t)=\alpha(1-t^2/T^2)$. We consider the family of surfaces $t\in [-T,T]\mapsto \Sigma_{h_{\alpha,T}(t)}$ which is locally parametrized by functions $(x,t)\mapsto \psi_{h_{\alpha,T}(t)}(x)$. Notice that it is a smooth family of smooth surfaces since due to Assumption~\eqref{listdef6.4.1} of Definition~\ref{defi-foliation1}. Also notice that the larger is $T$, the smaller are the derivatives of these functions with respect to $t$. By assumption, each function $x\mapsto \psi_\lambda(x)$ is non-characteristic and strongly pseudo-convex as a function independent of $t$. By compactness, there exists $T$ large enough such that $(x,t)\mapsto \psi_{h_{\alpha,T}(t)}(x)$ defines local surfaces which are non-characteristic and pseudo-convex for all $\alpha\in[0,\lambda_0]$ and $t\in[-T,T]$. The parameter $T$ is fixed in the remaining part of the proof and we may omit it in the notations.

Notice that, for any $\alpha$, the family of set $t\in[-T,T]\mapsto
\Sigma^+_{h_\alpha(t)}$ starts inside $\omega$ at $t=-T$ and finishes inside $\omega$ at $t=T$. Moreover, for any small $\alpha$, these sets always stay inside $\omega$ where $z$ vanishes. Assume that there exist $(x,t)$ and $\alpha\in[0,\lambda_0]$ such that $x\in
\Sigma^+_{h_\alpha(t)}$ and $z(x,t)\neq 0$. We set
\begin{equation}\label{eq-demo-global}
\alpha_0=\min\left\{\alpha\in (0,\lambda_0]\,,\,
\exists t\in[-T,T],\exists x\in\Sigma^+_{h_{\alpha}(t)}\text{ such that }z(x,t)\neq 0\right\}\,.
\end{equation}
By continuity, we know that $z(x,t)=0$ for all $x\in
\Sigma^+_{h_{\alpha_0}(t)}$. Moreover, there exists $t_0\in (-T,T)$ and $x_0\in\Sigma_{h_{\alpha_0}(t_0)}$ such that $z$ is not identically zero in any neighborhood of $(x_0,t_0)$. Indeed, otherwise, by compactness, we may extend the set where $z$ vanishes and contradict~\eqref{eq-demo-global}.

To conclude, it remains to use the local unique continuation property of Theorem~\ref{th-hormander} at $(x_0,t_0)$ with the time-space surface defined by $(x,t)\mapsto \psi_{h_{\alpha_0}(t)}(x)$. The continuation implies that $z$ vanishes near $(x_0,t_0)$ which contradicts the construction. Thus, $z(x,t)=0$ for all $t\in [-T,T]$ and $x\in \bigcup_\alpha \Sigma^+_{h_\alpha(t)}$. In particular $z(\cdot,\,t=0)\equiv 0$ in $\bigcup_{\alpha}\Sigma^+_{h_\alpha(0)}=\bigcup_{\lambda\leq\lambda_0}
\Sigma^+_\lambda$. Since these arguments hold for all $\lambda_0<1$ and since $\omega\cup_{\lambda\in [0,1)} \Sigma^+_\lambda$ is $\Omega$ up to a set of measure zero, we have that $z(\cdot,\,t=0)\equiv 0$ in $\Omega$. Well-posedness of the linear wave equation concludes that $z\equiv 0$ everywhere.
\end{proof}


Notice that, as it is stated, this unique continuation result needs an infinite time to be efficient, where Theorem~\ref{th-RZ} only need a finite explicit time. In fact, a careful look to the proof shows that a finite time is sufficient once we know that the family of surface is pseudo-convex and non-characteristic in a uniform way. However, such a bound of convexity is difficult to obtain in general cases and may be even impossible as for the example studied in Section~7.


A typical example of application is given in Figure~\ref{fig-foliation-1}: if $\Omega$ is a sphere and $\omega$ covers more than an hemisphere, then if $z$ is a global solution of a linear wave equation which vanishes in $\omega$ for all times, then $z\equiv 0$. Notice that, in this case, the family of surfaces is uniformly pseudo-convex and the unique continuation holds in fact in finite time even if $z$ is not a global in time solution.



\setcounter{subsection}{3}
\subsection{Proof of Theorem~\ref{th2}}

Theorem~\ref{th2} is then a direct consequence of the unique continuation results stated in this Section: Proposition~\ref{prop-UCP-implique-th} and Theorems~\ref{th-hormander-global} and~6.7 imply Theorem~\ref{th2}.



\begin{thebibliography}{MMMMMM}
\bibitem[AL14]{Ana-Leau} Anantharaman, Nalini; Léautaud, Matthieu Sharp polynomial decay rates for the damped wave equation on the torus, Anal. PDE, Volume 7 (2014) no. 1, pp. 159-214 \href{https://ahl.centre-mersenne.org/item/AHL_2020__3__1241_0/#Ana-Leau}{Publisher reference}.
\bibitem[Ana08]{Anant} Anantharaman, Nalini Entropy and the localization of eigenfunctions, Ann. Math., Volume 168 (2008) no. 2, pp. 435-475 \href{https://ahl.centre-mersenne.org/item/AHL_2020__3__1241_0/#Anant}{Publisher reference}.
\bibitem[BD08]{Batty-Duyckaerts} Batty, Charles J. K.; Duyckaerts, Thomas Non-uniform stability for bounded semi-groups on Banach spaces, J. Evol. Equ., Volume 8 (2008) no. 4, pp. 765-780 \href{https://ahl.centre-mersenne.org/item/AHL_2020__3__1241_0/#Batty-Duyckaerts}{Publisher reference}.
\bibitem[BD18]{BourgainDyatlov} Bourgain, Jean; Dyatlov, Semyon Spectral gaps without the pressure condition, Ann. of Math., Volume 187 (2018) no. 3, pp. 825-867 \href{https://ahl.centre-mersenne.org/item/AHL_2020__3__1241_0/#BourgainDyatlov}{Publisher reference}.
\bibitem[BL13]{Burq-Lebeau} Burq, Nicolas; Lebeau, Gilles Injections de Sobolev probabilistes et applications, Ann. Sci. Éc. Norm. Supér., Volume 46 (2013) no. 6, pp. 917-962 \href{https://ahl.centre-mersenne.org/item/AHL_2020__3__1241_0/#Burq-Lebeau}{Publisher reference}.
\bibitem[BLR92]{BLR} Bardos, Claude; Lebeau, Gilles; Rauch, Jeffrey Sharp sufficient conditions for the observation, control and stabilization of waves from the boundary, SIAM J. Control Optim., Volume 30 (1992) no. 5, pp. 1024-1065 \href{https://ahl.centre-mersenne.org/item/AHL_2020__3__1241_0/#BLR}{Publisher reference}.
\bibitem[BM04]{Burq-Zworski} Burq, Nicolas; Maciej, Zworski Geometric control in the presence of a black box, J. Am. Math. Soc., Volume 17 (2004) no. 2, pp. 443-471 \href{https://ahl.centre-mersenne.org/item/AHL_2020__3__1241_0/#Burq-Zworski}{Publisher reference}.
\bibitem[BT10]{Bor-Tom} Borichev, Alexander; Tomilov, Yury Optimal polynomial decay of functions and operator semigroups, Math. Ann., Volume 347 (2010) no. 2, pp. 455-478 \href{https://ahl.centre-mersenne.org/item/AHL_2020__3__1241_0/#Bor-Tom}{Publisher reference}.
\bibitem[Bur93]{Burq-memoire} Burq, Nicolas Contrôle de l’équation des plaques en présence d’obstacles strictement convexes, Mém. Soc. Math. France (N.S.), Volume 55 (1993), pp. 3-126 \href{https://ahl.centre-mersenne.org/item/AHL_2020__3__1241_0/#Burq-memoire}{Publisher reference}.
\bibitem[BV83]{Babin-Vishik} Babin, Anatolii V.; Vishik, Mark I. Regular attractors of semigroups and evolution equations, J. Math. Pures Appl., Volume 62 (1983) no. 4, pp. 441-491 \href{https://ahl.centre-mersenne.org/item/AHL_2020__3__1241_0/#Babin-Vishik}{Publisher reference}.
\bibitem[CKS + 10]{Colliander} Colliander, James E.; Keel, Markus; Staffilani, Gigliola; Takaoka, Hideo; Tao, Terence Transfer of energy to high frequencies in the cubic defocusing nonlinear Schrödinger equation, Invent. Math., Volume 181 (2010) no. 1, pp. 39-113 \href{https://ahl.centre-mersenne.org/item/AHL_2020__3__1241_0/#Colliander}{Publisher reference}.
\bibitem[CSVW14]{CSVW} Christianson, Hans; Schenck, Emmanuel; Vasy, András; Wunsch, Jared From resolvent estimates to damped waves, J. Anal. Math., Volume 1222 (2014), pp. 143-160 \href{https://ahl.centre-mersenne.org/item/AHL_2020__3__1241_0/#CSVW}{Publisher reference}.
\bibitem[Daf78]{Dafermos} Dafermos, Constantin M. Asymptotic behavior of solutions of evolution equations, Nonlinear Evolution Equations. Proc. Sympos., Univ. Wisconsin, Madison, Wis., 1977) (Publication of the Mathematics Research Center, the University of Wisconsin), Volume 40, Academic Press, 1978, pp. 103-123 \href{https://ahl.centre-mersenne.org/item/AHL_2020__3__1241_0/#Dafermos}{Publisher reference}.
\bibitem[Del01]{Dehman} Delham, Belhassen Stabilisation pour l’équation des ondes semi-linéaire, Asymptotic Anal., Volume 27 (2001) no. 2, pp. 171-181 \href{https://ahl.centre-mersenne.org/item/AHL_2020__3__1241_0/#Dehman}{Publisher reference}.
\bibitem[DLZ03]{DLZ} Delham, Belhassen; Lebeau, Gilles; Zuazua, Enrique Stabilization and control for the subcritical semilinear wave equation, Ann. Sci. Éc. Norm. Sup., Volume 36 (2003) no. 4, pp. 525-551 \href{https://ahl.centre-mersenne.org/item/AHL_2020__3__1241_0/#DLZ}{Publisher reference}.
\bibitem[DV12]{Vasy} Datchev, Kiril; Vasy, András Propagation through trapped sets and semiclassical resolvent estimates, Ann. Inst. Fourier, Volume 62 (2012) no. 6, pp. 2347-2377 (Addendum ibid. pp. 2379–2384) \href{https://ahl.centre-mersenne.org/item/AHL_2020__3__1241_0/#Vasy}{Publisher reference}.
\bibitem[Gea78]{Ge} Gearhart, Larry Spectral theory for contraction semigroups on Hilbert spaces, Trans. Am. Math. Soc., Volume 236 (1978), pp. 385-394 \href{https://ahl.centre-mersenne.org/item/AHL_2020__3__1241_0/#Ge}{Publisher reference}.
\bibitem[GG12]{Gerard-Grellier} Gérard, Patrick; Grellier, Sandrine Effective integrable dynamics for a certain nonlinear wave equation, Anal. PDE, Volume 5 (2012) no. 5, pp. 1139-1155 \href{https://ahl.centre-mersenne.org/item/AHL_2020__3__1241_0/#Gerard-Grellier}{Publisher reference}.
\bibitem[Hal88]{Hale-book} Hale, Jack K. Asymptotic behavior of dissipative systems, Mathematical Surveys and Monographs, 25, American Mathematical Society, 1988 \href{https://ahl.centre-mersenne.org/item/AHL_2020__3__1241_0/#Hale-book}{Publisher reference}.
\bibitem[Har85]{Haraux} Haraux, Alain Stabilization of trajectories for some weakly damped hyperbolic equations, J. Differ. Equ., Volume 59 (1985) no. 2, pp. 145-154 \href{https://ahl.centre-mersenne.org/item/AHL_2020__3__1241_0/#Haraux}{Publisher reference}.
\bibitem[HBTV15]{Hani} Hani, Zaher; Benoît, Pausader; Tzvetkov, Nikolay; Visciglia, Nicola Modified scattering for the cubic Schrödinger equation on product spaces and applications, Forum Math. Pi, Volume 3 (2015), e4, p. 63 \href{https://ahl.centre-mersenne.org/item/AHL_2020__3__1241_0/#Hani}{Publisher reference}.
\bibitem[HR03]{Hale-Raugel} Hale, Jack K.; Raugel, Geneviève Regularity, determining modes and Galerkin methods, J. Math. Pures Appl., Volume 82 (2003) no. 9, pp. 1075-1136 \href{https://ahl.centre-mersenne.org/item/AHL_2020__3__1241_0/#Hale-Raugel}{Publisher reference}.
\bibitem[Hua85]{Huang} Huang, Falun Characteristic conditions for exponential stability of linear dynamical systems in Hilbert spaces, Ann. Differ. Equations, Volume 1 (1985) no. 1, pp. 43-56 \href{https://ahl.centre-mersenne.org/item/AHL_2020__3__1241_0/#Huang}{Publisher reference}.
\bibitem[Hör85]{Hormander} Hörmander, Lars The analysis of linear partial differential operators. IV: Fourier integral operators, Grundlehren der Mathematischen Wissenschaften, 275, Springer, 1985 \href{https://ahl.centre-mersenne.org/item/AHL_2020__3__1241_0/#Hormander}{Publisher reference}.
\bibitem[Hör97]{Hormander_uniq_analytic} Hörmander, Lars On the uniqueness of the Cauchy problem under partial analyticity assumptions, Geometrical optics and related topics (Cortona, 1996 (Progress in Nonlinear Differential Equations and Their Applications), Volume 32, Birkhäuser (1997), pp. 179-219 \href{https://ahl.centre-mersenne.org/item/AHL_2020__3__1241_0/#Hormander_uniq_analytic}{Publisher reference}.
\bibitem[Ika82]{Ikawa1} Ikawa, Mitsuru Decay of solutions of the wave equation in the exterior of two convex obstacles, Osaka J. Math., Volume 19 (1982), pp. 459-509 \href{https://ahl.centre-mersenne.org/item/AHL_2020__3__1241_0/#Ikawa1}{Publisher reference}.
\bibitem[Ika88]{Ikawa2} Ikawa, Mitsuru Decay of solutions of the wave equation in the exterior of several convex bodies, Ann. Inst. Fourier, Volume 38 (1988) no. 2, pp. 113-146 \href{https://ahl.centre-mersenne.org/item/AHL_2020__3__1241_0/#Ikawa2}{Publisher reference}.
\bibitem[JL13]{RC} Joly, Romain; Laurent, Camille Stabilization for the semilinear wave equation with geometric control condition, Anal. PDE, Volume 6 (2013) no. 5, pp. 1089-1119 \href{https://ahl.centre-mersenne.org/item/AHL_2020__3__1241_0/#RC}{Publisher reference}.
\bibitem[Jol07]{RJ} Joly, Romain New examples of damped wave equations with gradient-like structure, Asymptotic Anal., Volume 53 (2007) no. 4, pp. 237-253 \href{https://ahl.centre-mersenne.org/item/AHL_2020__3__1241_0/#RJ}{Publisher reference}.
\bibitem[KK93]{K-K} Kazemi, Mohammad; Klibanov, Michael V. Stability estimates for ill-posed Cauchy problems involving hyperbolic equations and inequalities, Appl. Anal., Volume 50 (1993) no. 1-2, pp. 93-102 \href{https://ahl.centre-mersenne.org/item/AHL_2020__3__1241_0/#K-K}{Publisher reference}.
\bibitem[Leb93]{Lebeau} Lebeau, Gilles Equation des ondes amorties, Algebraic and Geometric Methods in Mathematical Physics: proceedings of the Kaciveli Summer School, Crimea, Ukraine, 1993 (Mathematical Physics Studies), Volume 19, Springer (1993), pp. 73-109 \href{https://ahl.centre-mersenne.org/item/AHL_2020__3__1241_0/#Lebeau}{Publisher reference}.
\bibitem[Lin20]{Lin} Lin, Jong Damped wave equations on compact hyperbolic surfaces, Commun. Math. Phys., Volume 373 (2020) no. 3, pp. 771-794 \href{https://ahl.centre-mersenne.org/item/AHL_2020__3__1241_0/#Lin}{Publisher reference}.
\bibitem[LL15]{Leautaud-Lerner} Léautaud, Matthieu; Lerner, Nicolas Sharp polynomial energy decay for locally undamped waves, Sémin. Équ. Dériv. Partielles (2014-2015), 21, p. 13 \href{https://ahl.centre-mersenne.org/item/AHL_2020__3__1241_0/#Leautaud-Lerner}{Publisher reference}.
\bibitem[LR85]{LR} Lerner, Nicolas; Robbiano, Luc Unicité de Cauchy pour des opérateurs de type principal, J. Anal. Math., Volume 44 (1985), pp. 32-66 \href{https://ahl.centre-mersenne.org/item/AHL_2020__3__1241_0/#LR}{Publisher reference}.
\bibitem[LR97]{Lebeau-Robbiano} Lebeau, Gilles; Robbiano, Luc Stabilisation de l’équation des ondes par le bord, Duke Math. J., Volume 86 (1997) no. 3, pp. 465-491 \href{https://ahl.centre-mersenne.org/item/AHL_2020__3__1241_0/#Lebeau-Robbiano}{Publisher reference}.
\bibitem[LTZ00]{Las-Tri-Zha} Lasiecka, Irena; Triggiani, Roberto; Zhang, Xu Nonconservative wave equations with unobserved Neumann B.C.: global uniqueness and observability in one shot, Differential geometric methods in the control of partial differential equations (Boulder, CO, 1999) (Contemporary Mathematics), Volume 268, American Mathematical Society, 2000, pp. 227-325 \href{https://ahl.centre-mersenne.org/item/AHL_2020__3__1241_0/#Las-Tri-Zha}{Publisher reference}.
\bibitem[Non18]{NonenmacherICM} Nonnenmacher, Stéphane Resonances in hyperbolic dynamics, Proceedings of the International Congress of Mathematicians (ICM 2018), Volume 2-4, World Scientific (2018), pp. 2495-2518 \href{https://ahl.centre-mersenne.org/item/AHL_2020__3__1241_0/#NonenmacherICM}{Publisher reference}.
\bibitem[Prü84]{Pruss} Prüss, Jan On the spectrum of C 0 -semigroups, Trans. Am. Math. Soc., Volume 284 (1984) no. 2, pp. 847-857 \href{https://ahl.centre-mersenne.org/item/AHL_2020__3__1241_0/#Pruss}{Publisher reference}.
\bibitem[Riv14]{Riv} Rivière, Gabriel Eigenmodes of the damped wave equation and small hyperbolic subsets, Ann. Inst. Fourier, Volume 64 (2014) no. 3, pp. 1229-1267 (With an appendix by Stéphane Nonnenmacher and Rivière) \href{https://ahl.centre-mersenne.org/item/AHL_2020__3__1241_0/#Riv}{Publisher reference}.
\bibitem[Rui92]{Ruiz} Ruiz, Alberto Unique continuation for weak solutions of the wave equation plus a potential, J. Math. Pures Appl., Volume 71 (1992) no. 5, pp. 455-467 \href{https://ahl.centre-mersenne.org/item/AHL_2020__3__1241_0/#Ruiz}{Publisher reference}.
\bibitem[RZ98]{Rob-Zui} Robbiano, Luc; Zuily, Claude Uniqueness in the Cauchy problem for operators with partially holomorphic coefficients, Invent. Math., Volume 131 (1998) no. 3, pp. 493-539 \href{https://ahl.centre-mersenne.org/item/AHL_2020__3__1241_0/#Rob-Zui}{Publisher reference}.
\bibitem[Sch10]{Schenckpressure} Schenck, Emmanuel Energy Decay for the Damped Wave Equation Under a Pressure Condition, Comm. Math. Phys., Volume 300 (2010) no. 2, pp. 375-410 \href{https://ahl.centre-mersenne.org/item/AHL_2020__3__1241_0/#Schenckpressure}{Publisher reference}.
\bibitem[Sch11]{Schenck} Schenck, Emmanuel Exponential stabilization without geometric control, Math. Res. Lett., Volume 18 (2011) no. 2, pp. 379-388 \href{https://ahl.centre-mersenne.org/item/AHL_2020__3__1241_0/#Schenck}{Publisher reference}.
\bibitem[SU13]{Stefanov-Uhlmann} Stefanov, Plamen; Ulhmann, Gunther Recovery of a source term or a speed with one measurement and applications, Trans. Am. Math. Soc., Volume 365 (2013) no. 11, pp. 5737-5758 \href{https://ahl.centre-mersenne.org/item/AHL_2020__3__1241_0/#Stefanov-Uhlmann}{Publisher reference}.
\bibitem[Tat95]{Tataru_uniq_analyticCPDE} Tataru, Daniel Unique continuation for solutions to PDE’s; between Hörmander’s theorem and Holmgren’s theorem, Commun. Partial Differ. Equations, Volume 20 (1995) no. 5-6, pp. 855-884 \href{https://ahl.centre-mersenne.org/item/AHL_2020__3__1241_0/#Tataru_uniq_analyticCPDE}{Publisher reference}.
\bibitem[Tat96]{Tataru} Tataru, Daniel Carleman estimates and unique continuation for solutions to boundary value problems, J. Math. Pures Appl., Volume 75 (1996) no. 4, pp. 367-408 \href{https://ahl.centre-mersenne.org/item/AHL_2020__3__1241_0/#Tataru}{Publisher reference}.
\bibitem[Tat99]{Tataru_uniq_analyticJMPA} Tataru, Daniel Unique continuation for operators with partially analytic coefficients, J. Math. Pures Appl., Volume 78 (1999) no. 5, pp. 505-521 \href{https://ahl.centre-mersenne.org/item/AHL_2020__3__1241_0/#Tataru_uniq_analyticJMPA}{Publisher reference}.
\bibitem[Zua90]{Zuazua} Zuazua, Enrique Exponential decay for semilinear wave equations with localized damping, Commun. Partial Differ. Equations, Volume 15 (1990) no. 2, pp. 205-235 \href{https://ahl.centre-mersenne.org/item/AHL_2020__3__1241_0/#Zuazua}{Publisher reference}.
\end{thebibliography}
\end{document}
